\documentclass[11pt]{article}
\usepackage{mathunal-base}
\mathunalfuente{serif}

\renewcommand{\examtitulo}{Cálculo en Varias Variables (1000006) --- Primer Examen Parcial}
\renewcommand{\examfecha}{Semestre 01-2013, 8 de Abril de 2013}

\begin{document}

\examheader

\vspace{4pt}
\noindent
\begin{minipage}[t]{0.62\linewidth}
\textbf{Duración: 1 hora y 50 minutos}

\vspace{8pt}
Nombre y carnet: \dotfill

\vspace{6pt}
Profesor y grupo: \dotfill
\end{minipage}

\vspace{8pt}
\noindent\textbf{Calificación:} 1a \dotfill; 1b \dotfill; 2a \dotfill; 2b \dotfill; 3a \dotfill; 3b \dotfill; 4 \dotfill. \textbf{Total:} \dotfill

\vspace{6pt}
\noindent\textbf{Instrucciones.} No está permitido el uso de calculadora ni de cualquier otro tipo de aparato electrónico, incluyendo teléfonos móviles, los cuales deben permanecer apagados durante el tiempo de duración del examen. Tampoco está permitido hacer preguntas a las personas encargadas de la vigilancia.

\vspace{4pt}
\noindent El examen consta de 4 puntos distribuidos en 4 páginas.

\begin{enumerate}
\item{} (25\%)
\begin{enumerate}
\item{} (12\%) Encuentre y dibuje el dominio de la función
\[
f(x,y) = \sqrt{4-x^2-y^2}+\ln(x^2+y^2-1).
\]
\item{} (13\%) Determine si existe
\[
\lim_{(x,y)\to(0,0)} \frac{2\operatorname{sen}(x^2+y^2)+x^4}{x^2+y^2}.
\]
\end{enumerate}

\vspace{5cm}

\newpage
\item{} (25\%) Sea $f:\mathbb{R}^2\longrightarrow\mathbb{R}$ la función definida por
\[
f(x,y) = \left\{
\begin{array}{ll}
\dfrac{x^3-2y^3}{x^2+y^2} & \text{si } (x,y)\neq(0,0),\\[6pt]
0 & \text{si } (x,y)=(0,0).
\end{array}
\right.
\]
\begin{enumerate}
\item{} (12\%) Encuentre las derivadas parciales $f_x(0,0)$ y $f_y(0,0)$.
\item{} (13\%) ¿Es $f$ diferenciable en $(0,0)$?
\end{enumerate}

\vspace{6cm}

\newpage
\item{} (25\%)
\begin{enumerate}
\item{} (12\%) ¿En cuál dirección la función $f(x,y)=1-x^2-y^2$ crece más rápidamente en $(1,1)$? Encuentre el valor de la máxima razón de cambio de $f$ en $(1,1)$.
\item{} (13\%) Encuentre los puntos en el hiperboloide $x^2+4y^2-z^2=4$ donde el plano tangente es paralelo al plano $2x+2y+z=5$.
\end{enumerate}

\vspace{6cm}

\newpage
\item{} (25\%) Sea $z=f(x,y)$ una función de clase $C^2$ y sean $x=r\cos\theta$, y $y=r\operatorname{sen}\theta$. Pruebe que
\[
z_{rr} = f_{xx}\cos^2\theta+2f_{xy}\cos\theta\operatorname{sen}\theta+f_{yy}\operatorname{sen}^2\theta.
\]

\vspace{6cm}
\end{enumerate}

\examfooter
\end{document}
